\documentclass[leqno]{article}
% Shared mark scheme preamble for resources.danieljsmith.org: the pack
% preamble, unchanged, then the scheme's own layer. A scheme's .tex is
% rendered from its YAML (scripts/README.md, Mark schemes), never edited.
\input{\CurrentFilePath/pack_preamble.tex}

% The page: equal side margins, so the table runs the full width of the
% header rule. The footers' right edge follows the text's.
\geometry{left=0.8in, right=0.8in}
\makeatletter
\renewcommand{\djs@footR}[1]{%
  \put(\dimexpr1in+\hoffset+\oddsidemargin+\textwidth\relax,
    -\dimexpr\paperheight-0.32in\relax){\makebox[0pt][r]{#1}}}
\makeatother

\usepackage[export]{adjustbox}
\usepackage{tabularray}
\definecolor{msGrid}{RGB}{176, 172, 164}
\definecolor{msTint}{RGB}{245, 243, 239}

% The header on every page, the front page included, and the booklet label
% under the contents-page title. \djsQuestionPageStyle stays empty.
\newcommand{\djsSchemeHeader}[1]{%
  \djsPackHeader[markscheme]{\textit{Mark scheme} - #1}%
  \renewcommand{\djsPackTitle}{#1}%
  \renewcommand{\djsBookletLabel}{Mark scheme}%
}

% The front page: the contents heading, spaced as tocloft spaces the packs'
% "Questions", then the contents table. A question's name and page number
% both link to its first page.
\newcommand{\msContentsHeading}{%
  \vskip\cftbeforetoctitleskip
  {\noindent\cfttoctitlefont\contentsname\par}%
  \vskip\cftaftertoctitleskip}
\NewTblrEnviron{mscontents}
\SetTblrInner[mscontents]{
  colspec = {Q[wd=26mm, l] X[l] Q[wd=16mm, c] Q[wd=14mm, r]},
  width = \linewidth,
  rowsep = 4pt, colsep = 6pt,
  hlines = {0.4pt, msGrid}, vlines = {0.4pt, msGrid},
  hline{2} = {2.2pt, scarlet},
  row{1} = {font=\small\bfseries},
  row{Z} = {font=\bfseries},
  column{2} = {font=\small},
}
\newcommand{\msContentsQuestion}[1]{\hyperref[q:#1]{Question #1}}
\newcommand{\msContentsPage}[1]{\hyperref[q:#1]{\pageref*{q:#1}}}
\newcommand{\msSep}{\hspace{0.9em}}

% The conventions pages' layout: headings in the documents' own bold, tables
% ruled and padded as the scheme's, roomy lists.
\newcommand{\msHead}[1]{%
  \par\vspace{18pt}%
  {\noindent\bfseries #1\par}%
  \vspace{6pt}}
\NewTblrEnviron{mstblr}
\SetTblrInner[mstblr]{rowsep = 6pt, colsep = 6pt,
  hlines = {0.4pt, msGrid}, vlines = {0.4pt, msGrid}}
\newenvironment{msList}
  {\begin{itemize}[leftmargin=1.2em, itemsep=5pt, topsep=3pt, parsep=0pt]}
  {\end{itemize}}

% A question's first page: the link target and the PDF bookmark, as
% \questionitem gives the packs.
\newcommand{\msQnum}{}
\newcommand{\msQuestion}[1]{%
  \clearpage
  \renewcommand{\msQnum}{#1}%
  \phantomsection\label{q:#1}%
  \addcontentsline{toc}{subsection}{\protect\djsTocEntry{#1}}}

% The scheme's table: Question 16mm, Scheme the rest, Marks 16mm, Notes
% 44mm; every cell ruled in msGrid; the scarlet rule under the head row. A
% long table has no caption or "continued" foot; a continuation page opens
% with a muted line above the repeated head row. The rules across a part's
% rows, the alternative's tint and the total row are the rendered table's.
\NewTblrTheme{ms}{%
  \DefTblrTemplate{firsthead}{default}{}%
  \DefTblrTemplate{middlehead,lasthead}{default}{%
    {\small\itshape\color{djsMuted}Question \msQnum\ continued\par}%
    \vspace{6pt}}%
  \DefTblrTemplate{firstfoot,middlefoot,lastfoot}{default}{}%
}
\NewTblrEnviron{msscheme}
\SetTblrOuter[msscheme]{long, theme=ms, entry=none, label=none}
\SetTblrInner[msscheme]{
  colspec = {Q[wd=16mm, l] X[l] Q[wd=16mm, c] Q[wd=44mm, l]},
  width = \linewidth,
  rowsep = 6pt, colsep = 6pt,
  rowhead = 1,
  vlines = {0.4pt, msGrid},
  hline{1} = {0.4pt, msGrid}, hline{2} = {2.2pt, scarlet},
  hline{Z} = {0.4pt, msGrid},
  column{1} = {font=\bfseries}, column{3} = {font=\bfseries},
  column{4} = {font=\small},
  row{1} = {rowsep=4pt, font=\small\bfseries},
}

% Two marks on a row: each note starts with its code.
\newcommand{\msNote}[2]{\textbf{#1}\enspace #2}
% A full-row figure: the worked solution's own picture, scaled uniformly and
% only down, to the table less the Question column. The picture is drawn
% with the line width it has in the solutions document, at its list depth
% there (the optional argument, 1 by default: a part's working), so one
% sized by \linewidth keeps its size; tabularray sets the cell's \linewidth
% to the cell's measured width, which would shrink it at every pass. The
% widths are measured once, in a solution box at the packs' text width
% (pack_preamble.tex's margins, 0.8in and 1in).
\newlength{\msSpanW}
\makeatletter
\newcommand{\ms@keepLW}[1]{%
  \expandafter\xdef\csname ms@solnLW@#1\endcsname{\the\linewidth}}
\AtBeginDocument{%
  \setlength{\msSpanW}{\dimexpr\textwidth-16mm-28pt\relax}%
  \setbox\z@\vbox{%
    \begin{minipage}{\dimexpr\paperwidth-1.8in\relax}
      \begin{djsSolution}
        \ms@keepLW{0}%
        \begin{enumerate}\item \ms@keepLW{1}%
        \begin{enumerate}\item \ms@keepLW{2}%
        \begin{enumerate}\item \ms@keepLW{3}%
        \end{enumerate}\end{enumerate}\end{enumerate}
      \end{djsSolution}
    \end{minipage}}}
\newcommand{\msFigure}[2][1]{%
  \adjustbox{max width=\msSpanW, max totalheight=0.6\textheight, valign=t}{%
    \hsize=\csname ms@solnLW@#1\endcsname
    \linewidth=\hsize
    #2}}
\makeatother
% A figure mark's feature, in the Scheme column.
\newcommand{\msFeature}[1]{\textit{#1}}
% An alternative's title row.
\newcommand{\msAlternative}[1]{\textit{Alternative:} #1}
% Drafts only: a queried mark's code, the query's why under its note, and a
% flag in its part's place.
\newcommand{\msQueried}[2]{#1~\(\to\) #2?}
\newcommand{\msQuery}[1]{{\itshape\color{djsMuted}Query: #1}}
\newcommand{\msFlag}[2]{\textbf{#1:}\enspace #2}


\djsSchemeHeader{Proof by Induction - Derivatives}
\djsSitePath{Core Pure}{Proof by Induction}

\djsLastUpdated{25 October 2026}
\begin{document}
\djsSetQuestionCount{8}
\djsAddMarks{1}{6}
\djsAddMarks{2}{6}
\djsAddMarks{3}{6}
\djsAddMarks{4}{7}
\djsAddMarks{5}{6}
\djsAddMarks{6}{11}
\djsAddMarks{7}{9}
\djsAddMarks{8}{7}
\vspace*{6pt}
\djsContentsTitle
\vspace{30pt}
\msContentsHeading
\begin{mscontents}{}
Question & Marks by part & Marks & Page \\
\msContentsQuestion{1} & & 6 & \msContentsPage{1} \\
\msContentsQuestion{2} & & 6 & \msContentsPage{2} \\
\msContentsQuestion{3} & & 6 & \msContentsPage{3} \\
\msContentsQuestion{4} & & 7 & \msContentsPage{4} \\
\msContentsQuestion{5} & & 6 & \msContentsPage{5} \\
\msContentsQuestion{6} & (a)\,3\msSep (b)\,6\msSep (c)\,2 & 11 & \msContentsPage{6} \\
\msContentsQuestion{7} & (a)\,6\msSep (b)\,3 & 9 & \msContentsPage{7} \\
\msContentsQuestion{8} & & 7 & \msContentsPage{8} \\
\SetCell[c=2]{r} Total & & 58 & \\
\end{mscontents}
\newpage
% The conventions pages, shared by every pack's mark scheme: one source of
% truth, \input after the front page, with no strand or pack logic. The keys
% on the first side, the rules on the second. The scheme layer defines
% \msHead, msList and mstblr.
\vspace*{6pt}%
{\noindent\cfttoctitlefont Reading the mark scheme\par}%
\vspace{10pt}%
These rules apply throughout, unless a note says otherwise.

\msHead{Marks}
Every mark is worth one, so \textbf{M1 A1} on a row is two marks. Only marks
coded ft follow through.\par
\vspace{10pt}%
\begin{mstblr}{
  colspec = {Q[wd=15mm, l] X[l]},
  width = \linewidth,
  hline{2} = {2.2pt, scarlet},
  column{1} = {font=\bfseries},
  row{1} = {font=\small\bfseries, rowsep=4pt},
}
Code & Meaning \\
M1 & Method: a correct method, applied to this question. A slip doesn't cost it if the method is clear. A formula must be used, not just quoted. \\
dM1 & Dependent method: as M1, but only if the method mark before it was earned, unless its note names another. \\
A1 & Accuracy: a correct answer or step. It needs the method mark it follows. \\
A1ft & As A1, following through from an earlier answer that the note names. \\
A1~AG & The last mark of a part whose answer is printed in the question: complete working, with no errors. \\
B1 & A correct value, fact or statement. It needs no method mark. \\
B1ft & As B1, following through from an earlier answer that the note names. \\
E1 & A written explanation, interpretation or conclusion. Concluding a proof by induction needs every earlier mark. \\
\end{mstblr}

\msHead{Abbreviations}
ft and AG qualify a code in the Marks column. The others appear in the
notes.\par
\vspace{10pt}%
\begin{mstblr}{
  colspec = {Q[wd=17mm, l] X[l] Q[wd=15mm, l] X[l]},
  width = \linewidth,
  column{1} = {font=\bfseries}, column{3} = {font=\bfseries},
}
ft & follow through & AG & answer given in the question \\
cao & correct answer only, with no follow-through & cso & correct solution only: no errors in the part \\
oe & or an equivalent form & awrt & any answer that rounds to this value \\
soi & seen, or implied by later working & isw & ignore working after a correct answer \\
SC & special case: credit outside the main method & their & the student's own earlier value \\
condone & a slip that doesn't cost the mark & [\;\,] & working the mark doesn't need \\
s.f. & significant figures & d.p. & decimal places \\
\end{mstblr}

\newpage
\msHead{Reading the scheme}
\begin{msList}
\item Each row is a line of working, with the marks it earns beside it.
  Where a row carries two marks, each note starts with its code.
\item \textbf{[\textit{n}]} ends each part with its total.
\item A shaded \textit{Alternative} is a second method for the same total.
  Mark an answer by the method it follows.
\item Any other valid method earns the equivalent marks, unless the question
  names a method or says \textit{hence}.
\item A sketch earns one B1 for each feature named below its figure.
\end{msList}

\msHead{Accuracy and answers}
\begin{msList}
\item Unless a question says otherwise, give answers to 3 significant
  figures, and angles in degrees to 1 decimal place. Angles in radians take
  3 significant figures.
\item A final answer with more figures is accepted if it rounds to the
  scheme's value; one with fewer is not. An answer spoiled by rounding too
  early loses its accuracy mark.
\item Where a question asks for an exact answer, a decimal loses the final
  accuracy mark. Otherwise either is accepted.
\item Wrong or missing units cost nothing unless the question asks for them.
\item A correct answer with no working earns full marks, unless the question
  needs the working: show that, prove, hence, a method it names, or an exact
  answer. A wrong answer with no working earns nothing. Calculator answers
  are treated the same way.
\item A correct answer, once reached, isn't spoiled by later working (isw).
\end{msList}

\msHead{Value of $\boldsymbol{g}$}
\begin{msList}
\item Take $g = 9.8\,\mathrm{m\,s^{-2}}$ unless a question states a value,
  which is then the only one accepted. Answers that use $g = 9.8$ may be
  given to 2 or 3 significant figures.
\item Where a question doesn't state $g$, an answer worked with $g = 10$
  throughout earns the same marks, keeping any value printed in the
  question. The worked solutions list the final answers with $g = 10$.
\end{msList}


\msQuestion{1}
\begin{msscheme}{
  hline{3-6} = {2-Z}{0.4pt, msGrid},
  hline{7} = {0.4pt, msGrid},
  row{7} = {font=\bfseries},
}
Question & Scheme & Marks & Notes \\
1 & {\(\frac{\mathrm{d}y}{\mathrm{d}x} = (ax^2+2x)e^{ax}\)\\[4pt] \(\frac{\mathrm{d}^2y}{\mathrm{d}x^2} = (a^2x^2+4ax+2)e^{ax}\)\\[4pt] \(n=2\): \(a^0(a^2x^2+4ax+2)e^{ax}\), so true for \(n=2\)} & B1 & {Second derivative found by differentiating and shown equal to the formula at \(n=2\). Not for substituting \(n=2\) into the formula alone.} \\
& {Assume true for \(n=k\):\\[4pt] \(\frac{\mathrm{d}^{k+1}y}{\mathrm{d}x^{k+1}} = \frac{\mathrm{d}}{\mathrm{d}x}\left[a^{k-2}\left(a^2x^2+2kax+k(k-1)\right)e^{ax}\right]\)} & M1 & {States the result assumed for \(n=k\) and takes its derivative as the \((k+1)\)th derivative.} \\
& {\(\phantom{=}a^{k-2}\left[(2a^2x+2ka)e^{ax}\right.\)\\[4pt] \(\phantom{=a^{k-2}}\left.+\left(a^2x^2+2kax+k(k-1)\right)ae^{ax}\right]\)} & M1 A1 & {\msNote{M1}{Applies the product rule to their \(k\)th derivative: two terms, the quadratic differentiated and \(e^{ax}\) differentiated to \(ae^{ax}\). Condone coefficient slips.}\\ \msNote{A1}{Any correct unsimplified form, \(a^{k-2}\) factored out or multiplied in.}} \\
& {\(= a^{k-1}\left(a^2x^2+2(k+1)ax+k(k+1)\right)e^{ax}\)\\[4pt] \(= a^{(k+1)-2}\left(a^2x^2+2(k+1)ax\right.\)\\[4pt] \(\phantom{=a^{(k+1)-2}}\left.+(k+1)\left((k+1)-1\right)\right)e^{ax}\)} & A1 & {cso. Reaches \(a^{k-1}(\ldots)e^{ax}\) and shows it is the formula with \(n=k+1\), e.g. \(k(k+1)\) written as \((k+1)((k+1)-1)\) or the target stated.} \\
& {True for \(n=2\); true for \(n=k\) \(\Rightarrow\) true for \(n=k+1\)\\[4pt] \(\Rightarrow\) true for all integers \(n \ge 2\)} & E1 & {Concludes naming true for \(n=2\), true for \(n=k\) implies true for \(n=k+1\), and hence true for all integers \(n \ge 2\).} \\
\nopagebreak \SetCell[c=2]{r} Total for Question 1 & & 6 & \\
\end{msscheme}

\msQuestion{2}
\begin{msscheme}{
  hline{3-6} = {2-Z}{0.4pt, msGrid},
  hline{7} = {0.4pt, msGrid},
  row{7} = {font=\bfseries},
}
Question & Scheme & Marks & Notes \\
2 & {\(\frac{\mathrm{d}}{\mathrm{d}x}(x^3e^{-x}) = 3x^2e^{-x} - x^3e^{-x}\)\\[4pt] \(\phantom{\frac{\mathrm{d}}{\mathrm{d}x}(x^3e^{-x})} = -e^{-x}(x^3-3x^2)\)\\[4pt] \(n=1\): \((-1)^1e^{-x}(x^3-3x^2+0-0)\), so true} & B1 & {Needs both the derivative and the formula evaluated at \(n=1\), shown equal.} \\
& {Assume true for \(n=k\):\\[4pt] \(\frac{\mathrm{d}^{k+1}}{\mathrm{d}x^{k+1}}(x^3e^{-x})\)\\[4pt] \(\phantom{A} = \frac{\mathrm{d}}{\mathrm{d}x}\left[(-1)^ke^{-x}(x^3-3kx^2+\ldots)\right]\)} & M1 & {States the assumption for \(n=k\) and differentiates its expression to obtain the \((k+1)\)th derivative. Not for assuming \(n=k+1\).} \\
& {\(= (-1)^k\bigl[-e^{-x}(x^3-3kx^2+3k(k-1)x\)\\[4pt] \(\phantom{=} -k(k-1)(k-2)) + e^{-x}(3x^2-6kx+3k(k-1))\bigr]\)} & M1 A1 & {\msNote{M1}{Uses the product rule on \(e^{-x}\) times their cubic: two terms, one with the cubic's derivative. Condone sign slips.}\\ \msNote{A1}{Correct, unsimplified.}} \\
& {\(= (-1)^{k+1}e^{-x}\bigl(x^3-3(k+1)x^2\)\\[4pt] \(\phantom{=} +3(k+1)kx-(k+1)k(k-1)\bigr)\)} & A1 & {cso. Each coefficient collected to the \(n=k+1\) form with \((-1)^{k+1}\); constant \(-(k+1)k(k-1)\) oe accepted.} \\
& {True for \(n=1\); true for \(n=k\) \(\Rightarrow\) true for \(n=k+1\); so true for all positive integers \(n\)} & E1 & {States the base case, the implication from \(n=k\) to \(n=k+1\), and truth for all positive integers \(n\).} \\
\nopagebreak \SetCell[c=2]{r} Total for Question 2 & & 6 & \\
\end{msscheme}

\msQuestion{3}
\begin{msscheme}{
  hline{3-6} = {2-Z}{0.4pt, msGrid},
  hline{7} = {0.4pt, msGrid},
  row{7} = {font=\bfseries},
}
Question & Scheme & Marks & Notes \\
3 & {\(n=1\): \(\frac{\mathrm{d}y}{\mathrm{d}x} = e^{-2x}(-2\cosh 3x + 3\sinh 3x)\)\\[4pt] \(\frac{1+(-5)}{2} = -2,\ \frac{1-(-5)}{2} = 3\), so true for \(n=1\)} & B1 & {Correct derivative and both coefficients of the formula evaluated at \(n=1\), seen to agree.} \\
& {Assume \(\frac{\mathrm{d}^k y}{\mathrm{d}x^k} = e^{-2x}\Big(\tfrac{1+(-5)^k}{2}\cosh 3x\)\\[4pt] \(\phantom{=}\ + \tfrac{1-(-5)^k}{2}\sinh 3x\Big)\)\\[4pt] \(\frac{\mathrm{d}^{k+1}y}{\mathrm{d}x^{k+1}} = \frac{\mathrm{d}}{\mathrm{d}x}\left(\frac{\mathrm{d}^k y}{\mathrm{d}x^k}\right)\)} & M1 & {States the result assumed for \(n=k\) and sets out to differentiate it to obtain the \((k+1)\)th derivative.} \\
& {\(\frac{\mathrm{d}^{k+1}y}{\mathrm{d}x^{k+1}} = -2e^{-2x}\Big(\tfrac{1+(-5)^k}{2}\cosh 3x\)\\[4pt] \(\phantom{=}\ + \tfrac{1-(-5)^k}{2}\sinh 3x\Big)\)\\[4pt] \(\phantom{=}\ + e^{-2x}\Big(\tfrac{3(1+(-5)^k)}{2}\sinh 3x\)\\[4pt] \(\phantom{=}\ + \tfrac{3(1-(-5)^k)}{2}\cosh 3x\Big)\)} & M1 A1 & {\msNote{M1}{Applies the product rule to the \(n=k\) expression: two terms, the second with \(\cosh 3x\) and \(\sinh 3x\) interchanged. Condone coefficient and sign slips.}\\ \msNote{A1}{Correct, unsimplified.}} \\
& {\(= e^{-2x}\Big(\tfrac{1-5(-5)^k}{2}\cosh 3x + \tfrac{1+5(-5)^k}{2}\sinh 3x\Big)\)\\[4pt] \(= e^{-2x}\Big(\tfrac{1+(-5)^{k+1}}{2}\cosh 3x\)\\[4pt] \(\phantom{=}\ + \tfrac{1-(-5)^{k+1}}{2}\sinh 3x\Big)\)} & A1 & {cso. Both coefficients collected and reached in the \((-5)^{k+1}\) form without error; not for jumping to the printed form.} \\
& {True for \(n=1\); true for \(n=k \Rightarrow\) true for \(n=k+1\)\\[4pt] \(\Rightarrow\) true for all \(n\in\mathbb{N}\)} & E1 & {States all three: true for \(n=1\); true for \(n=k\) implies true for \(n=k+1\); hence true for all \(n\in\mathbb{N}\) oe.} \\
\nopagebreak \SetCell[c=2]{r} Total for Question 3 & & 6 & \\
\end{msscheme}

\msQuestion{4}
\begin{msscheme}{
  hline{3-7} = {2-Z}{0.4pt, msGrid},
  hline{8} = {0.4pt, msGrid},
  row{8} = {font=\bfseries},
}
Question & Scheme & Marks & Notes \\
4 & {\(\frac{\mathrm{d}y}{\mathrm{d}x} = 2x\cosh x + x^2\sinh x\)\\[4pt] \(\frac{\mathrm{d}^2y}{\mathrm{d}x^2} = x^2\cosh x + 4x\sinh x + 2\cosh x\)\\[4pt] Compares with the formula at \(n = 1\)} & B1 & {Both derivatives correct and matched to the formula at \(n = 1\). Not for substituting \(n = 1\) into the formula alone.} \\
& {Assumes \(\frac{\mathrm{d}^{2k}y}{\mathrm{d}x^{2k}} = x^2\cosh x + 4kx\sinh x\)\\[4pt] \(\phantom{\frac{\mathrm{d}^{2k}y}{\mathrm{d}x^{2k}}} + 2k(2k-1)\cosh x\)\\[4pt] \(\frac{\mathrm{d}^{2k+2}y}{\mathrm{d}x^{2k+2}} = \frac{\mathrm{d}^2}{\mathrm{d}x^2}\left(\frac{\mathrm{d}^{2k}y}{\mathrm{d}x^{2k}}\right)\)} & M1 & {States the result for \(n = k\) and takes the \((2k+2)\)th derivative as its second derivative; implied by differentiating it.} \\
& {\(\frac{\mathrm{d}^{2k+1}y}{\mathrm{d}x^{2k+1}} = 2x\cosh x + x^2\sinh x\)\\[4pt] \(\phantom{=} + 4k(\sinh x + x\cosh x) + 2k(2k-1)\sinh x\)} & M1 & {Differentiates their \(n = k\) expression, with the product rule on at least one product and \(\cosh x \to \sinh x\). Condone coefficient slips.} \\
& {\(\frac{\mathrm{d}^{2k+2}y}{\mathrm{d}x^{2k+2}} = 2x\sinh x + x^2\cosh x\)\\[4pt] \(\phantom{=} + (4k+2)(\cosh x + x\sinh x) + 2k(2k+1)\cosh x\)} & M1 A1 & {\msNote{M1}{Differentiates their \((2k+1)\)th derivative in the same way. Condone coefficient slips.}\\ \msNote{A1}{Correct, in any unsimplified form.}} \\
& {\(4k^2 + 6k + 2 = 2(k+1)\left(2(k+1)-1\right)\)\\[4pt] \(\Rightarrow \frac{\mathrm{d}^{2k+2}y}{\mathrm{d}x^{2k+2}} = x^2\cosh x + 4(k+1)x\sinh x\)\\[4pt] \(\phantom{\Rightarrow} + 2(k+1)\left(2(k+1)-1\right)\cosh x\)} & A1 & {cso. Collects to the formula with \(n = k+1\), the \(\cosh x\) coefficient shown as \(2(k+1)(2k+1)\) or better.} \\
& {True for \(n = 1\); true for \(n = k\) \(\Rightarrow\) true for \(n = k+1\);\\[4pt] so true for all integers \(n \ge 1\)} & E1 & {States all three: true for \(n = 1\), true for \(n = k\) implies true for \(n = k+1\), hence true for all \(n \ge 1\).} \\
\nopagebreak \SetCell[c=2]{r} Total for Question 4 & & 7 & \\
\end{msscheme}

\msQuestion{5}
\begin{msscheme}{
  hline{3-6} = {2-Z}{0.4pt, msGrid},
  hline{7} = {0.4pt, msGrid},
  row{7} = {font=\bfseries},
}
Question & Scheme & Marks & Notes \\
5 & {\(\frac{\mathrm{d}y}{\mathrm{d}x} = 2a(ax+1)\ln(ax+1) + a(ax+1)\)\\[4pt] \(\frac{\mathrm{d}^2y}{\mathrm{d}x^2} = a^2\bigl(2\ln(ax+1)+3\bigr)\)\\[4pt] \(\frac{\mathrm{d}^3y}{\mathrm{d}x^3} = \frac{2a^3}{ax+1}\)\\[4pt] \(n=3\): \(2(-1)^0\frac{0!\,a^3}{(ax+1)^1} = \frac{2a^3}{ax+1}\)} & B1 & {Three correct derivatives, any equivalent form, and the formula at \(n=3\) shown equal to the third derivative.} \\
& {Assume \(\frac{\mathrm{d}^ky}{\mathrm{d}x^k} = 2(-1)^{k-3}\frac{(k-3)!\,a^k}{(ax+1)^{k-2}}\)\\[4pt] \(\frac{\mathrm{d}^{k+1}y}{\mathrm{d}x^{k+1}} = \frac{\mathrm{d}}{\mathrm{d}x}\left(2(-1)^{k-3}(k-3)!\,a^k(ax+1)^{-(k-2)}\right)\)} & M1 & {States the \(n=k\) assumption and attempts to differentiate it to find the \((k+1)\)th derivative.} \\
& {\(\phantom{\frac{\mathrm{d}^{k+1}y}{\mathrm{d}x^{k+1}}} = 2(-1)^{k-3}(k-3)!\,a^k\)\\[4pt] \(\phantom{\frac{\mathrm{d}^{k+1}y}{\mathrm{d}x^{k+1}} =} \times\bigl(-(k-2)\bigr)(ax+1)^{-(k-1)}\times a\)} & M1 A1 & {\msNote{M1}{Differentiates to the form \(c(ax+1)^{-(k-1)}\). Condone sign slips and a missing factor \(a\).}\\ \msNote{A1}{Correct, unsimplified.}} \\
& {\(\phantom{\frac{\mathrm{d}^{k+1}y}{\mathrm{d}x^{k+1}}} = 2(-1)^{k-2}\frac{(k-2)!\,a^{k+1}}{(ax+1)^{k-1}}\)} & A1 & {cso. Needs \((-1)^{k-2}\) from the sign and \((k-2)(k-3)! = (k-2)!\), shown or clearly implied.} \\
& {True for \(n=3\); true for \(n=k\) \(\Rightarrow\) true for \(n=k+1\);\\[4pt] so true for all integers \(n \ge 3\)} & E1 & {States all three points: the base case, the implication from \(n=k\) to \(n=k+1\), and truth for all integers \(n \ge 3\).} \\
\nopagebreak \SetCell[c=2]{r} Total for Question 5 & & 6 & \\
\end{msscheme}

\msQuestion{6}
\begin{msscheme}{
  hline{3-5,7-11,13-14} = {2-Z}{0.4pt, msGrid},
  hline{6,12,15} = {0.4pt, msGrid},
  row{15} = {font=\bfseries},
}
Question & Scheme & Marks & Notes \\
6(a) & {\(\frac{\mathrm{d}y}{\mathrm{d}t} = a\ln(1+at) + a\)\\[4pt] \(\frac{\mathrm{d}^2y}{\mathrm{d}t^2} = a^2(1+at)^{-1}\), \(\frac{\mathrm{d}^3y}{\mathrm{d}t^3} = \ldots\)} & M1 & {Differentiates repeatedly, reaching at least the second and third derivatives in the form \(ka^m(1+at)^{-p}\). Condone sign and coefficient slips.} \\
& {\(\frac{\mathrm{d}^2y}{\mathrm{d}t^2} = \frac{a^2}{1+at}\), \(\frac{\mathrm{d}^3y}{\mathrm{d}t^3} = -\frac{a^3}{(1+at)^2}\)\\[4pt] \(\frac{\mathrm{d}^4y}{\mathrm{d}t^4} = \frac{2a^4}{(1+at)^3}\), \(\frac{\mathrm{d}^5y}{\mathrm{d}t^5} = -\frac{6a^5}{(1+at)^4}\)} & A1 & {All four correct, any equivalent form.} \\
& {\(\frac{\mathrm{d}^ny}{\mathrm{d}t^n} = \frac{(-1)^n(n-2)!\,a^n}{(1+at)^{n-1}}\)} & A1 & {oe, such as \((-1)^{n-2}\) or \((1+at)^{1-n}\) in the numerator. Must be in terms of \(n\), \(a\) and \(t\).} \\
\nopagebreak & & [3] & \\
6(b) & {\(n=2\): \(\frac{(-1)^2\,0!\,a^2}{(1+at)^1} = \frac{a^2}{1+at} = f''(t)\)} & B1 & {Shows the formula with \(n=2\) gives \(\frac{a^2}{1+at}\) and that this matches \(f''(t)\).} \\
& {Assume \(f^{(k)}(t) = \frac{(-1)^k(k-2)!\,a^k}{(1+at)^{k-1}}\)\\[4pt] \(f^{(k+1)}(t)\)\\[4pt] \(\phantom{A} = \frac{\mathrm{d}}{\mathrm{d}t}\left[(-1)^k(k-2)!\,a^k(1+at)^{-(k-1)}\right]\)} & M1 & {States the formula assumed for \(n=k\) and sets out to differentiate it to obtain \(f^{(k+1)}(t)\).} \\
& {\(\phantom{A} = (-1)^k(k-2)!\,a^k\)\\[4pt] \(\phantom{A = } \times\left(-(k-1)\right)a(1+at)^{-k}\)} & M1 A1 & {\msNote{M1}{Differentiates the assumed result, reducing the power of \((1+at)\) by 1 with a multiplying factor. Condone sign and coefficient slips.}\\ \msNote{A1}{Correct derivative, unsimplified, including the factor \(a\).}} \\
& {\(\phantom{A} = \frac{(-1)^{k+1}(k-1)!\,a^{k+1}}{(1+at)^{k}}\)} & A1 & {cso. Combines \((k-1)(k-2)!\) as \((k-1)!\) and the signs as \((-1)^{k+1}\), giving the formula with \(n=k+1\).} \\
& {True for \(n=2\); true for \(n=k\)\\[4pt] \(\Rightarrow\) true for \(n=k+1\);\\[4pt] so true for all integers \(n\ge 2\)} & E1 & {Concludes, naming true for \(n=2\), true for \(n=k\) implies true for \(n=k+1\), and hence true for all integers \(n\ge 2\).} \\
\nopagebreak & & [6] & \\
6(c) & {\(f^{(7)}\left(\tfrac53\right) = \frac{(-1)^7\,5!\times 3^7}{\left(1+3\times\frac53\right)^6} = \ldots\)} & M1 & {Substitutes \(n=7\), \(a=3\) and \(t=\frac53\) into their formula. Not for the sixth derivative.} \\
& {\(-\frac{45}{8}\)} & A1 & {oe, or awrt \(-5.63\).} \\
\nopagebreak & & [2] & \\
\nopagebreak \SetCell[c=2]{r} Total for Question 6 & & 11 & \\
\end{msscheme}

\msQuestion{7}
\begin{msscheme}{
  hline{3-7,9-11} = {2-Z}{0.4pt, msGrid},
  hline{8,12} = {0.4pt, msGrid},
  row{12} = {font=\bfseries},
}
Question & Scheme & Marks & Notes \\
7(a) & {\(f'(x)=-\frac12(1+x)^{-\frac32}\)\\[4pt] \(\frac{(-1)^1\,2!}{2^2\,1!\,(1+x)^{\frac32}}=-\frac12(1+x)^{-\frac32}\)} & B1 & {Both the derivative and the formula at \(n=1\), shown equal.} \\
& {Assume \(f^{(k)}(x)=\frac{(-1)^k(2k)!}{2^{2k}k!(1+x)^{k+\frac12}}\)\\[4pt] \(f^{(k+1)}(x)=\frac{(-1)^k(2k)!}{2^{2k}k!}\times\frac{\mathrm{d}}{\mathrm{d}x}(1+x)^{-k-\frac12}\)} & M1 & {States the formula assumed for \(n=k\) and sets out to differentiate it for \(f^{(k+1)}(x)\). Not for assuming the \(n=k+1\) result.} \\
& {\(\phantom{f^{(k+1)}(x)}=\frac{(-1)^k(2k)!}{2^{2k}k!}\left(-k-\frac12\right)\)\\[4pt] \(\phantom{f^{(k+1)}(x)=}\times(1+x)^{-k-\frac32}\)} & M1 A1 & {\msNote{M1}{Differentiates their \(f^{(k)}(x)\) to the form \(c(1+x)^{-k-\frac32}\). Condone sign and coefficient slips.}\\ \msNote{A1}{Any unsimplified form.}} \\
& {\(\phantom{f^{(k+1)}(x)}=\frac{(-1)^{k+1}(2k+2)(2k+1)(2k)!}{2^{2k+2}(k+1)k!(1+x)^{k+\frac32}}\)\\[4pt] \(\phantom{f^{(k+1)}(x)}=\frac{(-1)^{k+1}(2k+2)!}{2^{2k+2}(k+1)!(1+x)^{k+\frac32}}\)} & A1 & {cso: the \(n=k+1\) formula reached with \((2k+2)!\) and \((k+1)!\) formed explicitly; any order of steps. Jumping to the result scores A0.} \\
& {True for \(n=1\); true for \(n=k\Rightarrow\) true for \(n=k+1\)\\[4pt] \(\Rightarrow\) true for all \(n\ge1\)} & E1 & {States all three points; condone equivalent wording such as ``for all positive integers \(n\)''.} \\
\nopagebreak & & [6] & \\
7(b) & {\(f(0)=1,\quad f^{(n)}(0)=\frac{(-1)^n(2n)!}{2^{2n}n!}\)} & B1 & {Both, with \(f^{(n)}(0)\) from part (a).} \\
& {\(\frac{x^n}{n!}f^{(n)}(0)=\frac{x^n}{n!}\times\frac{(-1)^n(2n)!}{2^{2n}n!}\)} & M1 & {Substitutes their \(f^{(n)}(0)\) into \(\frac{x^n}{n!}f^{(n)}(0)\), the Maclaurin general term quoted or implied.} \\
& {\(\phantom{\frac{x^n}{n!}f^{(n)}(0)}=(-1)^n\frac{(2n)!}{2^{2n}(n!)^2}x^n\)\\[4pt] \((1+x)^{-\frac12}=1-\frac{x}{2}+\frac{3x^2}{8}-\cdots\)} & A1~AG & {Correct substituted general term, combining the two \(n!\) to \((n!)^2\), and the printed series stated without error.} \\
\nopagebreak & & [3] & \\
\nopagebreak \SetCell[c=2]{r} Total for Question 7 & & 9 & \\
\end{msscheme}

\msQuestion{8}
\begin{msscheme}{
  hline{3-8} = {2-Z}{0.4pt, msGrid},
  hline{9} = {0.4pt, msGrid},
  row{9} = {font=\bfseries},
}
Question & Scheme & Marks & Notes \\
8 & {\(\frac{\mathrm{d}^2}{\mathrm{d}x^2}(x^2f) = 2f+4xf'+x^2f''\)\\[4pt] \(\phantom{\frac{\mathrm{d}^2}{\mathrm{d}x^2}(x^2f)} = 2f+4xf'-x^2f\)\\[4pt] \(\phantom{\frac{\mathrm{d}^2}{\mathrm{d}x^2}(x^2f)} = -(x^2f-4xf'-2f)\)} & B1 & {Differentiates twice, uses \(f''=-f\) and shows agreement with the RHS at \(n=1\). Not for substituting \(n=1\) into the RHS alone.} \\
& {Assume \(\frac{\mathrm{d}^{2k}}{\mathrm{d}x^{2k}}(x^2f)\)\\[4pt] \(\phantom{A} = (-1)^k(x^2f-4kxf'-2k(2k-1)f)\)\\[4pt] \(\frac{\mathrm{d}^{2k+2}}{\mathrm{d}x^{2k+2}}(x^2f)\)\\[4pt] \(\phantom{A} = (-1)^k\left[\frac{\mathrm{d}^2}{\mathrm{d}x^2}(x^2f)-4k\frac{\mathrm{d}^2}{\mathrm{d}x^2}(xf')\right.\)\\[4pt] \(\phantom{A = (-1)^k} \left.-2k(2k-1)f''\right]\)} & M1 & {Assumes the result for \(n=k\) and writes the \((2k+2)\)th derivative as the second derivative of the assumed expression.} \\
& {\(\frac{\mathrm{d}^2}{\mathrm{d}x^2}(x^2f) = 2f+4xf'+x^2f''\)\\[4pt] \(\frac{\mathrm{d}^2}{\mathrm{d}x^2}(xf') = 2f''+xf'''\)} & M1 & {Finds the second derivatives of \(x^2f\) (or quotes the base case) and \(xf'\) by the product rule; whole-expression differentiation also scores. Condone slips.} \\
& {\(f''=-f,\ f'''=-f' \Rightarrow \frac{\mathrm{d}^{2k+2}}{\mathrm{d}x^{2k+2}}(x^2f) = \ldots\)} & M1 & {Replaces \(f''\) by \(-f\) and \(f'''\) by \(-f'\) throughout and collects terms in \(x^2f\), \(xf'\) and \(f\). Condone sign and coefficient slips.} \\
& {\(\phantom{A} = (-1)^k(-x^2f+4(k+1)xf'\)\\[4pt] \(\phantom{A = (-1)^k(} +(2+8k+2k(2k-1))f)\)} & A1 & {Any correct unsimplified form in \(f\) and \(f'\) only.} \\
& {\(\phantom{A} = (-1)^{k+1}(x^2f-4(k+1)xf'\)\\[4pt] \(\phantom{A = (-1)^{k+1}(} -2(k+1)(2k+1)f)\)} & A1 & {cso, with the coefficient of \(f\) seen as \(2(k+1)(2k+1)\) or \(2(k+1)(2(k+1)-1)\).} \\
& {True for \(n=1\); true for \(n=k \Rightarrow\) true for \(n=k+1\);\\[4pt] so true for all positive integers \(n\)} & E1 & {States true for \(n=1\), true for \(n=k\) implies true for \(n=k+1\), hence true for all positive integers \(n\).} \\
\nopagebreak \SetCell[c=2]{r} Total for Question 8 & & 7 & \\
\end{msscheme}

\end{document}
